\section{A third model: d1.2-ultra}\label{sec:third}

After both studies had been analysed, a third model of the same kind was released by a third maker: d1.2-ultra.
It answers the same kind of request as the two models of the analysis, so it could be put to the same test
without a new study. Nothing in this section was fixed in advance, and no verdict changes with it.

\paragraph{What was done.} d1.2-ultra rated the 5{,}000 candidate excerpts of the pool once, at 250 words and on
the whole question bank, with the request that \jev{} and \drex{} had received, and the 160 excerpts of the
second study three more times. Its scores were then compared with people's ratings exactly as the other two
models' had been. Its dimension scores are computed as those of \drex{} are: the scorer of the analysis, which
was fitted on \jev{}'s ratings, applied to the model's ratings standardised within the model. On dimensions
\jev{} therefore has the home side; on single questions the three models stand on the same footing.

\begin{figure}[!htb]\centering
\begin{tikzpicture}[fig]
\mkdot{cjev}{4.7}{1.66}
\node[leg,anchor=base west] at (4.93,1.57) {jev-1.14.0};
\mkdia{cdrex}{6.77}{1.66}
\node[leg,anchor=base west] at (7,1.57) {drex-max};
\fill[caccent!85!black] (8.476,1.56) rectangle (8.676,1.76);
\node[leg,anchor=base west] at (8.806,1.57) {d1.2-ultra};
\mkiv{cpeople}{10.496}{10.796}{1.66}
\node[leg,anchor=base west] at (10.876,1.57) {95\% interval};
\draw[grid] (0,0) -- (15.65,0);
\draw[grid] (0,-0.68) -- (15.65,-0.68);
\draw[grid] (0,-1.36) -- (15.65,-1.36);
\draw[grid] (0,-2.04) -- (15.65,-2.04);
\draw[grid] (0,-2.72) -- (15.65,-2.72);
\draw[grid] (0,-3.4) -- (15.65,-3.4);
\draw[grid] (0,-4.08) -- (15.65,-4.08);
\draw[grid] (0,-4.76) -- (15.65,-4.76);
\draw[grid] (0,-5.44) -- (15.65,-5.44);
\node[ttl,anchor=base west] at (4.6,0.95) {Dimension score and composite};
\draw[grid] (4.6,0) -- (4.6,-5.44);
\node[tick,anchor=north] at (4.6,-5.56) {0};
\draw[grid] (5.6,0) -- (5.6,-5.44);
\node[tick,anchor=north] at (5.6,-5.56) {0.2};
\draw[grid] (6.6,0) -- (6.6,-5.44);
\node[tick,anchor=north] at (6.6,-5.56) {0.4};
\draw[grid] (7.6,0) -- (7.6,-5.44);
\node[tick,anchor=north] at (7.6,-5.56) {0.6};
\draw[grid] (8.6,0) -- (8.6,-5.44);
\node[tick,anchor=north] at (8.6,-5.56) {0.8};
\draw[grid] (9.6,0) -- (9.6,-5.44);
\node[tick,anchor=north] at (9.6,-5.56) {1.0};
\draw[axis] (4.6,-5.44) -- (9.6,-5.44);
\node[mut,anchor=north] at (7.1,-6.02) {Correlation with the composite mean};
\draw[ref] (7.1,0.12) -- (7.1,-5.44);
\node[ttl,anchor=base west] at (10.05,0.95) {d1.2-ultra minus};
\node[mut,anchor=base west] at (10.05,0.59) {jev-1.14.0};
\node[ttl,anchor=base west] at (13.05,0.95) {};
\node[mut,anchor=base west] at (13.05,0.59) {drex-max};
\node[lbl,anchor=base east] at (4.4,-0.43) {F1 Direct, organized argument};
\mkiv{cjev}{7.694}{8.326}{-0.17}
\mkdot{cjev}{8.016}{-0.17}
\mkiv{cdrex}{7.852}{8.491}{-0.34}
\mkdia{cdrex}{8.177}{-0.34}
\mkiv{caccent!85!black}{7.565}{8.298}{-0.51}
\path[draw=white,line width=0.6pt,fill=caccent!85!black] (7.953,-0.51) ++(-2.1pt,-2.1pt) rectangle ++(4.2pt,4.2pt);
\node[lbl,anchor=base west] at (10.05,-0.43) {\textminus 0.01 {\scriptsize[\textminus 0.07, +0.04]}};
\node[lbl,anchor=base west] at (13.05,-0.43) {\textminus 0.04 {\scriptsize[\textminus 0.09, \textminus 0.01]}};
\node[lbl,anchor=base east] at (4.4,-1.11) {F2 Positive tone vs disappointment};
\mkiv{cjev}{8.253}{8.797}{-0.85}
\mkdot{cjev}{8.549}{-0.85}
\mkiv{cdrex}{8.03}{8.669}{-1.02}
\mkdia{cdrex}{8.378}{-1.02}
\mkiv{caccent!85!black}{8.154}{8.725}{-1.19}
\path[draw=white,line width=0.6pt,fill=caccent!85!black] (8.452,-1.19) ++(-2.1pt,-2.1pt) rectangle ++(4.2pt,4.2pt);
\node[lbl,anchor=base west] at (10.05,-1.11) {\textminus 0.02 {\scriptsize[\textminus 0.06, +0.02]}};
\node[lbl,anchor=base west] at (13.05,-1.11) {+0.01 {\scriptsize[\textminus 0.01, +0.05]}};
\node[lbl,anchor=base east] at (4.4,-1.79) {F3 Self-disclosure, self-criticism};
\mkiv{cjev}{8.267}{8.866}{-1.53}
\mkdot{cjev}{8.591}{-1.53}
\mkiv{cdrex}{8.359}{8.842}{-1.7}
\mkdia{cdrex}{8.632}{-1.7}
\mkiv{caccent!85!black}{8.405}{8.881}{-1.87}
\path[draw=white,line width=0.6pt,fill=caccent!85!black] (8.666,-1.87) ++(-2.1pt,-2.1pt) rectangle ++(4.2pt,4.2pt);
\node[lbl,anchor=base west] at (10.05,-1.79) {+0.01 {\scriptsize[\textminus 0.02, +0.06]}};
\node[lbl,anchor=base west] at (13.05,-1.79) {+0.01 {\scriptsize[\textminus 0.02, +0.03]}};
\node[lbl,anchor=base east] at (4.4,-2.47) {F4 Care and regard for others};
\mkiv{cjev}{6.136}{7.941}{-2.21}
\mkdot{cjev}{7.109}{-2.21}
\mkiv{cdrex}{6.114}{7.936}{-2.38}
\mkdia{cdrex}{7.104}{-2.38}
\mkiv{caccent!85!black}{6.56}{8.06}{-2.55}
\path[draw=white,line width=0.6pt,fill=caccent!85!black] (7.385,-2.55) ++(-2.1pt,-2.1pt) rectangle ++(4.2pt,4.2pt);
\node[lbl,anchor=base west] at (10.05,-2.47) {+0.06 {\scriptsize[\textminus 0.02, +0.16]}};
\node[lbl,anchor=base west] at (13.05,-2.47) {+0.06 {\scriptsize[+0.01, +0.12]}};
\node[lbl,anchor=base east] at (4.4,-3.15) {F5 Narrative vs rebuttal};
\mkiv{cjev}{7.523}{8.473}{-2.89}
\mkdot{cjev}{8.084}{-2.89}
\mkiv{cdrex}{7.191}{8.307}{-3.06}
\mkdia{cdrex}{7.823}{-3.06}
\mkiv{caccent!85!black}{7.247}{8.422}{-3.23}
\path[draw=white,line width=0.6pt,fill=caccent!85!black] (7.933,-3.23) ++(-2.1pt,-2.1pt) rectangle ++(4.2pt,4.2pt);
\node[lbl,anchor=base west] at (10.05,-3.15) {\textminus 0.03 {\scriptsize[\textminus 0.09, +0.01]}};
\node[lbl,anchor=base west] at (13.05,-3.15) {+0.02 {\scriptsize[\textminus 0.02, +0.06]}};
\node[lbl,anchor=base east] at (4.4,-3.83) {F6 Energetic and angry vs calm};
\mkiv{cjev}{8.451}{8.895}{-3.57}
\mkdot{cjev}{8.693}{-3.57}
\mkiv{cdrex}{8.055}{8.606}{-3.74}
\mkdia{cdrex}{8.35}{-3.74}
\mkiv{caccent!85!black}{8.12}{8.659}{-3.91}
\path[draw=white,line width=0.6pt,fill=caccent!85!black] (8.413,-3.91) ++(-2.1pt,-2.1pt) rectangle ++(4.2pt,4.2pt);
\node[lbl,anchor=base west] at (10.05,-3.83) {\textminus 0.06 {\scriptsize[\textminus 0.09, \textminus 0.02]}};
\node[lbl,anchor=base west] at (13.05,-3.83) {+0.01 {\scriptsize[\textminus 0.02, +0.04]}};
\node[lbl,anchor=base east] at (4.4,-4.51) {F7 Humor, imaginative language};
\mkiv{cjev}{6.831}{7.907}{-4.25}
\mkdot{cjev}{7.399}{-4.25}
\mkiv{cdrex}{6.371}{7.6}{-4.42}
\mkdia{cdrex}{7.066}{-4.42}
\mkiv{caccent!85!black}{6.536}{7.69}{-4.59}
\path[draw=white,line width=0.6pt,fill=caccent!85!black] (7.164,-4.59) ++(-2.1pt,-2.1pt) rectangle ++(4.2pt,4.2pt);
\node[lbl,anchor=base west] at (10.05,-4.51) {\textminus 0.05 {\scriptsize[\textminus 0.12, +0.02]}};
\node[lbl,anchor=base west] at (13.05,-4.51) {+0.02 {\scriptsize[\textminus 0.04, +0.08]}};
\node[lbl,anchor=base east] at (4.4,-5.19) {F8 Hedging vs certainty};
\mkiv{cjev}{6.549}{7.727}{-4.93}
\mkdot{cjev}{7.151}{-4.93}
\mkiv{cdrex}{6.972}{8.003}{-5.1}
\mkdia{cdrex}{7.493}{-5.1}
\mkiv{caccent!85!black}{6.961}{7.978}{-5.27}
\path[draw=white,line width=0.6pt,fill=caccent!85!black] (7.483,-5.27) ++(-2.1pt,-2.1pt) rectangle ++(4.2pt,4.2pt);
\node[lbl,anchor=base west] at (10.05,-5.19) {+0.07 {\scriptsize[+0.02, +0.11]}};
\node[lbl,anchor=base west] at (13.05,-5.19) {\textminus 0.00 {\scriptsize[\textminus 0.05, +0.05]}};
\end{tikzpicture}%

\caption{Second study: each model's dimension score against people's composite, with 95\% intervals from
resampling recordings. Right: d1.2-ultra's correlation minus that of each earlier model, with the interval of
the difference.}\label{fig:s3dims}
\end{figure}

\paragraph{On the eight dimensions the three models are level.} Over the eight dimensions d1.2-ultra correlates
with people's composite at 0.666, \jev{} at 0.670 and \drex{} at 0.656 (Figure~\ref{fig:s3dims};
Table~\ref{tab:s3dims}). The difference from \jev{} is $-0.004$ ($-0.02$ to $+0.02$) and from \drex{} $+0.011$
($-0.00$ to $+0.03$). On F4, where the two earlier models sit at the line of 0.50, d1.2-ultra is at 0.56 (0.39
to 0.69): further from the line, with the same wide interval. Applied with each model alone, the rule gives the
same verdict for every dimension.

\begin{figure}[!htb]\centering
\begin{tikzpicture}[fig]
\mkdot{cjev}{3}{1.66}
\node[leg,anchor=base west] at (3.23,1.57) {jev-1.14.0};
\mkdia{cdrex}{5.07}{1.66}
\node[leg,anchor=base west] at (5.3,1.57) {drex-max};
\fill[caccent!85!black] (6.776,1.56) rectangle (6.976,1.76);
\node[leg,anchor=base west] at (7.106,1.57) {d1.2-ultra};
\mkiv{cpeople}{8.796}{9.096}{1.66}
\node[leg,anchor=base west] at (9.176,1.57) {95\% interval};
\draw[grid] (0,0) -- (15.35,0);
\draw[grid] (0,-0.62) -- (15.35,-0.62);
\draw[grid] (0,-1.24) -- (15.35,-1.24);
\draw[grid] (0,-1.86) -- (15.35,-1.86);
\draw[grid] (0,-2.48) -- (15.35,-2.48);
\draw[grid] (0,-3.1) -- (15.35,-3.1);
\draw[grid] (0,-3.72) -- (15.35,-3.72);
\draw[grid] (0,-4.34) -- (15.35,-4.34);
\draw[grid] (0,-4.96) -- (15.35,-4.96);
\draw[grid] (0,-5.58) -- (15.35,-5.58);
\draw[grid] (0,-6.2) -- (15.35,-6.2);
\draw[grid] (0,-6.82) -- (15.35,-6.82);
\draw[grid] (0,-7.44) -- (15.35,-7.44);
\draw[grid] (0,-8.06) -- (15.35,-8.06);
\draw[grid] (0,-8.68) -- (15.35,-8.68);
\draw[grid] (0,-9.3) -- (15.35,-9.3);
\draw[grid] (0,-9.92) -- (15.35,-9.92);
\node[ttl,anchor=base west] at (2.9,0.95) {Question score with the mean of five raters};
\draw[grid] (2.9,0) -- (2.9,-9.92);
\node[tick,anchor=north] at (2.9,-10.04) {0};
\draw[grid] (4.18,0) -- (4.18,-9.92);
\node[tick,anchor=north] at (4.18,-10.04) {0.2};
\draw[grid] (5.46,0) -- (5.46,-9.92);
\node[tick,anchor=north] at (5.46,-10.04) {0.4};
\draw[grid] (6.74,0) -- (6.74,-9.92);
\node[tick,anchor=north] at (6.74,-10.04) {0.6};
\draw[grid] (8.02,0) -- (8.02,-9.92);
\node[tick,anchor=north] at (8.02,-10.04) {0.8};
\draw[grid] (9.3,0) -- (9.3,-9.92);
\node[tick,anchor=north] at (9.3,-10.04) {1.0};
\draw[axis] (2.9,-9.92) -- (9.3,-9.92);
\node[mut,anchor=north] at (6.1,-10.5) {Correlation with the mean of five raters};
\draw[ref] (6.1,0.12) -- (6.1,-9.92);
\node[ttl,anchor=base west] at (9.75,0.95) {d1.2-ultra minus};
\node[mut,anchor=base west] at (9.75,0.59) {jev-1.14.0};
\node[ttl,anchor=base west] at (12.75,0.95) {};
\node[mut,anchor=base west] at (12.75,0.59) {drex-max};
\node[lbl,anchor=base east] at (2.7,-0.4) {Persuasion};
\mkiv{cjev}{6.247}{7.497}{-0.14}
\mkdot{cjev}{6.912}{-0.14}
\mkiv{cdrex}{5.985}{7.355}{-0.31}
\mkdia{cdrex}{6.646}{-0.31}
\mkiv{caccent!85!black}{6.1}{7.433}{-0.48}
\path[draw=white,line width=0.6pt,fill=caccent!85!black] (6.75,-0.48) ++(-2.1pt,-2.1pt) rectangle ++(4.2pt,4.2pt);
\node[lbl,anchor=base west] at (9.75,-0.4) {\textminus 0.03 {\scriptsize[\textminus 0.09, +0.04]}};
\node[lbl,anchor=base west] at (12.75,-0.4) {+0.02 {\scriptsize[\textminus 0.05, +0.07]}};
\node[lbl,anchor=base east] at (2.7,-1.02) {Persistence};
\mkiv{cjev}{6.485}{7.57}{-0.76}
\mkdot{cjev}{7.046}{-0.76}
\mkiv{cdrex}{5.636}{7.251}{-0.93}
\mkdia{cdrex}{6.527}{-0.93}
\mkiv{caccent!85!black}{6.257}{7.433}{-1.1}
\path[draw=white,line width=0.6pt,fill=caccent!85!black] (6.889,-1.1) ++(-2.1pt,-2.1pt) rectangle ++(4.2pt,4.2pt);
\node[lbl,anchor=base west] at (9.75,-1.02) {\textminus 0.02 {\scriptsize[\textminus 0.11, +0.05]}};
\node[lbl,anchor=base west] at (12.75,-1.02) {+0.06 {\scriptsize[\textminus 0.03, +0.16]}};
\node[lbl,anchor=base east] at (2.7,-1.64) {Positive tone};
\mkiv{cjev}{7.98}{8.576}{-1.38}
\mkdot{cjev}{8.313}{-1.38}
\mkiv{cdrex}{7.421}{8.219}{-1.55}
\mkdia{cdrex}{7.858}{-1.55}
\mkiv{caccent!85!black}{7.67}{8.425}{-1.72}
\path[draw=white,line width=0.6pt,fill=caccent!85!black] (8.098,-1.72) ++(-2.1pt,-2.1pt) rectangle ++(4.2pt,4.2pt);
\node[lbl,anchor=base west] at (9.75,-1.64) {\textminus 0.03 {\scriptsize[\textminus 0.06, \textminus 0.00]}};
\node[lbl,anchor=base west] at (12.75,-1.64) {+0.04 {\scriptsize[+0.01, +0.07]}};
\node[lbl,anchor=base east] at (2.7,-2.26) {Disappointment};
\mkiv{cjev}{7.036}{8.034}{-2}
\mkdot{cjev}{7.579}{-2}
\mkiv{cdrex}{6.883}{7.873}{-2.17}
\mkdia{cdrex}{7.402}{-2.17}
\mkiv{caccent!85!black}{7.168}{8.06}{-2.34}
\path[draw=white,line width=0.6pt,fill=caccent!85!black] (7.643,-2.34) ++(-2.1pt,-2.1pt) rectangle ++(4.2pt,4.2pt);
\node[lbl,anchor=base west] at (9.75,-2.26) {+0.01 {\scriptsize[\textminus 0.02, +0.04]}};
\node[lbl,anchor=base west] at (12.75,-2.26) {+0.04 {\scriptsize[\textminus 0.00, +0.08]}};
\node[lbl,anchor=base east] at (2.7,-2.88) {Self-disclosure};
\mkiv{cjev}{7.538}{8.263}{-2.62}
\mkdot{cjev}{7.931}{-2.62}
\mkiv{cdrex}{8.244}{8.674}{-2.79}
\mkdia{cdrex}{8.474}{-2.79}
\mkiv{caccent!85!black}{8.026}{8.536}{-2.96}
\path[draw=white,line width=0.6pt,fill=caccent!85!black] (8.303,-2.96) ++(-2.1pt,-2.1pt) rectangle ++(4.2pt,4.2pt);
\node[lbl,anchor=base west] at (9.75,-2.88) {+0.06 {\scriptsize[+0.02, +0.10]}};
\node[lbl,anchor=base west] at (12.75,-2.88) {\textminus 0.03 {\scriptsize[\textminus 0.05, \textminus 0.00]}};
\node[lbl,anchor=base east] at (2.7,-3.5) {Self-critical};
\mkiv{cjev}{7.37}{8.294}{-3.24}
\mkdot{cjev}{7.917}{-3.24}
\mkiv{cdrex}{6.878}{8.056}{-3.41}
\mkdia{cdrex}{7.58}{-3.41}
\mkiv{caccent!85!black}{7.54}{8.414}{-3.58}
\path[draw=white,line width=0.6pt,fill=caccent!85!black] (8.053,-3.58) ++(-2.1pt,-2.1pt) rectangle ++(4.2pt,4.2pt);
\node[lbl,anchor=base west] at (9.75,-3.5) {+0.02 {\scriptsize[\textminus 0.02, +0.06]}};
\node[lbl,anchor=base west] at (12.75,-3.5) {+0.07 {\scriptsize[+0.04, +0.12]}};
\node[lbl,anchor=base east] at (2.7,-4.12) {Compassion};
\mkiv{cjev}{5.353}{7.54}{-3.86}
\mkdot{cjev}{6.598}{-3.86}
\mkiv{cdrex}{5.618}{7.799}{-4.03}
\mkdia{cdrex}{6.859}{-4.03}
\mkiv{caccent!85!black}{6.14}{7.824}{-4.2}
\path[draw=white,line width=0.6pt,fill=caccent!85!black] (7.108,-4.2) ++(-2.1pt,-2.1pt) rectangle ++(4.2pt,4.2pt);
\node[lbl,anchor=base west] at (9.75,-4.12) {+0.08 {\scriptsize[+0.01, +0.20]}};
\node[lbl,anchor=base west] at (12.75,-4.12) {+0.04 {\scriptsize[\textminus 0.02, +0.12]}};
\node[lbl,anchor=base east] at (2.7,-4.74) {Respect\textsuperscript{n}};
\mkiv{cjev}{4.575}{6.395}{-4.48}
\mkdot{cjev}{5.582}{-4.48}
\mkiv{cdrex}{4.313}{6.187}{-4.65}
\mkdia{cdrex}{5.272}{-4.65}
\mkiv{caccent!85!black}{5.166}{6.777}{-4.82}
\path[draw=white,line width=0.6pt,fill=caccent!85!black] (6.06,-4.82) ++(-2.1pt,-2.1pt) rectangle ++(4.2pt,4.2pt);
\node[lbl,anchor=base west] at (9.75,-4.74) {+0.07 {\scriptsize[\textminus 0.03, +0.17]}};
\node[lbl,anchor=base west] at (12.75,-4.74) {+0.12 {\scriptsize[+0.03, +0.22]}};
\node[lbl,anchor=base east] at (2.7,-5.36) {Narrative};
\mkiv{cjev}{5.029}{7.017}{-5.1}
\mkdot{cjev}{6.163}{-5.1}
\mkiv{cdrex}{5.301}{7.357}{-5.27}
\mkdia{cdrex}{6.531}{-5.27}
\mkiv{caccent!85!black}{5.586}{7.489}{-5.44}
\path[draw=white,line width=0.6pt,fill=caccent!85!black] (6.712,-5.44) ++(-2.1pt,-2.1pt) rectangle ++(4.2pt,4.2pt);
\node[lbl,anchor=base west] at (9.75,-5.36) {+0.09 {\scriptsize[+0.01, +0.16]}};
\node[lbl,anchor=base west] at (12.75,-5.36) {+0.03 {\scriptsize[\textminus 0.01, +0.08]}};
\node[lbl,anchor=base east] at (2.7,-5.98) {Rebuttal\textsuperscript{n}};
\mkiv{cjev}{6.093}{7.486}{-5.72}
\mkdot{cjev}{6.888}{-5.72}
\mkiv{cdrex}{5.682}{7.35}{-5.89}
\mkdia{cdrex}{6.567}{-5.89}
\mkiv{caccent!85!black}{6.268}{7.529}{-6.06}
\path[draw=white,line width=0.6pt,fill=caccent!85!black] (6.98,-6.06) ++(-2.1pt,-2.1pt) rectangle ++(4.2pt,4.2pt);
\node[lbl,anchor=base west] at (9.75,-5.98) {+0.01 {\scriptsize[\textminus 0.04, +0.08]}};
\node[lbl,anchor=base west] at (12.75,-5.98) {+0.06 {\scriptsize[+0.01, +0.13]}};
\node[lbl,anchor=base east] at (2.7,-6.6) {Energetic};
\mkiv{cjev}{6.713}{7.791}{-6.34}
\mkdot{cjev}{7.315}{-6.34}
\mkiv{cdrex}{5.972}{7.284}{-6.51}
\mkdia{cdrex}{6.685}{-6.51}
\mkiv{caccent!85!black}{6.225}{7.472}{-6.68}
\path[draw=white,line width=0.6pt,fill=caccent!85!black] (6.911,-6.68) ++(-2.1pt,-2.1pt) rectangle ++(4.2pt,4.2pt);
\node[lbl,anchor=base west] at (9.75,-6.6) {\textminus 0.06 {\scriptsize[\textminus 0.11, \textminus 0.01]}};
\node[lbl,anchor=base west] at (12.75,-6.6) {+0.04 {\scriptsize[\textminus 0.01, +0.09]}};
\node[lbl,anchor=base east] at (2.7,-7.22) {Calm};
\mkiv{cjev}{6.946}{7.815}{-6.96}
\mkdot{cjev}{7.416}{-6.96}
\mkiv{cdrex}{6.875}{7.697}{-7.13}
\mkdia{cdrex}{7.324}{-7.13}
\mkiv{caccent!85!black}{7.049}{7.972}{-7.3}
\path[draw=white,line width=0.6pt,fill=caccent!85!black] (7.538,-7.3) ++(-2.1pt,-2.1pt) rectangle ++(4.2pt,4.2pt);
\node[lbl,anchor=base west] at (9.75,-7.22) {+0.02 {\scriptsize[\textminus 0.02, +0.05]}};
\node[lbl,anchor=base west] at (12.75,-7.22) {+0.03 {\scriptsize[\textminus 0.01, +0.08]}};
\node[lbl,anchor=base east] at (2.7,-7.84) {Humor};
\mkiv{cjev}{6.383}{7.565}{-7.58}
\mkdot{cjev}{7.037}{-7.58}
\mkiv{cdrex}{6.172}{7.427}{-7.75}
\mkdia{cdrex}{6.867}{-7.75}
\mkiv{caccent!85!black}{6.378}{7.585}{-7.92}
\path[draw=white,line width=0.6pt,fill=caccent!85!black] (7.059,-7.92) ++(-2.1pt,-2.1pt) rectangle ++(4.2pt,4.2pt);
\node[lbl,anchor=base west] at (9.75,-7.84) {+0.00 {\scriptsize[\textminus 0.06, +0.07]}};
\node[lbl,anchor=base west] at (12.75,-7.84) {+0.03 {\scriptsize[\textminus 0.01, +0.07]}};
\node[lbl,anchor=base east] at (2.7,-8.46) {Imaginative\textsuperscript{n}};
\mkiv{cjev}{4.701}{6.582}{-8.2}
\mkdot{cjev}{5.701}{-8.2}
\mkiv{cdrex}{4.462}{6.527}{-8.37}
\mkdia{cdrex}{5.609}{-8.37}
\mkiv{caccent!85!black}{4.819}{6.633}{-8.54}
\path[draw=white,line width=0.6pt,fill=caccent!85!black] (5.788,-8.54) ++(-2.1pt,-2.1pt) rectangle ++(4.2pt,4.2pt);
\node[lbl,anchor=base west] at (9.75,-8.46) {+0.01 {\scriptsize[\textminus 0.06, +0.08]}};
\node[lbl,anchor=base west] at (12.75,-8.46) {+0.03 {\scriptsize[\textminus 0.03, +0.10]}};
\node[lbl,anchor=base east] at (2.7,-9.08) {Hedging};
\mkiv{cjev}{5.719}{7.089}{-8.82}
\mkdot{cjev}{6.425}{-8.82}
\mkiv{cdrex}{5.799}{7.064}{-8.99}
\mkdia{cdrex}{6.512}{-8.99}
\mkiv{caccent!85!black}{5.973}{7.122}{-9.16}
\path[draw=white,line width=0.6pt,fill=caccent!85!black] (6.597,-9.16) ++(-2.1pt,-2.1pt) rectangle ++(4.2pt,4.2pt);
\node[lbl,anchor=base west] at (9.75,-9.08) {+0.03 {\scriptsize[\textminus 0.03, +0.08]}};
\node[lbl,anchor=base west] at (12.75,-9.08) {+0.01 {\scriptsize[\textminus 0.03, +0.06]}};
\node[lbl,anchor=base east] at (2.7,-9.7) {Knowledge limits\textsuperscript{n}};
\mkiv{cjev}{2.9}{4.719}{-9.44}
\mkdot{cjev}{3.648}{-9.44}
\mkiv{cdrex}{3.096}{5.219}{-9.61}
\mkdia{cdrex}{4.257}{-9.61}
\mkiv{caccent!85!black}{3.09}{5.235}{-9.78}
\path[draw=white,line width=0.6pt,fill=caccent!85!black] (4.215,-9.78) ++(-2.1pt,-2.1pt) rectangle ++(4.2pt,4.2pt);
\node[lbl,anchor=base west] at (9.75,-9.7) {+0.09 {\scriptsize[+0.02, +0.16]}};
\node[lbl,anchor=base west] at (12.75,-9.7) {\textminus 0.01 {\scriptsize[\textminus 0.05, +0.04]}};
\end{tikzpicture}%

\caption{Second study: each model's score on a question against the mean of five raters, with 95\% intervals.
Right: d1.2-ultra minus each earlier model. \textsuperscript{n}: not rated in the first
study.}\label{fig:s3questions}
\end{figure}

\paragraph{On single questions d1.2-ultra is slightly ahead.} Over the sixteen questions of the second study it
correlates with the mean of five raters at 0.628, against 0.606 for \jev{} and 0.591 for \drex{}: $+0.022$
($+0.00$ to $+0.04$) and $+0.036$ ($+0.02$ to $+0.05$) (Figure~\ref{fig:s3questions}). It is the best of the
three on ten of the sixteen. Over the twelve questions of the first study the figures are 0.687, 0.674 and
0.669 (Table~\ref{tab:s3first}). The lead over \drex{} is clear in both studies. The lead over \jev{} is small:
its interval starts at zero in the second study and includes zero in the first.

\paragraph{Its levels are closer to people's.} Over the sixteen questions d1.2-ultra's mean score is 0.18 scale
points away from people's mean rating on average, against 0.37 for \jev{} and 0.33 for \drex{}. Its largest
distance is $+0.41$ on self-disclosure.

\paragraph{It shares the other models' blind spots.} With rater noise taken out, d1.2-ultra's misses correlate
with \jev{}'s at 0.64 to 0.92 over the eight dimensions. On F4 it shows the same difference between streams and
videos at the same score ($-0.46$ standard deviations, $-0.79$ to $-0.13$). Four calls of the same text agree as
closely as for the other two models (Table~\ref{tab:repeat}).

\paragraph{What follows.} A third model from a third maker reads these excerpts as people do to the same degree
as the first two, misses where they miss, and differs from them mainly in its levels. What the two studies found
is therefore not a property of one model: where five readers agree, each of three models follows them, and
where readers do not agree, a newer model does not help.
